\subsection{Anodyne extensions}

The homotopy theory of simplicial sets revolves around fibrations.
Originally this was developed in terms of the extension condition, which results in tedious combinatorial manipulations.
Fortunately with a more modern viewpoint we can suppress most of these combinatorial arguments.

\bdefn

    A class $M$ of monomorphisms of $\sSet$ is called {\emphcolor saturated} if the following conditions are satisfied.
    \benum
        \item $M$ contains all isomorphisms.
        \item $M$ is closed under pushouts in the sense that if

        \commsq{A&C\cr B&B\cup_AC}
            {}{}{i}{i_*}

        is a pushout square and $i\in M$, then $i_*\in M$ as well.

        \item Each retract of an element of $M$ is in $M$.
        That is, given the following commutative diagram

        \centerline{\drawdiagram{
            $A'$&$A$&$A'$\cr
            $B'$&$B$&$B'$\cr
        }{
            \diagarrow{from={1,1}, to={2,1}, text={i'}, x distance=.25cm}
            \diagarrow{from={1,2}, to={2,2}, text={i}, x distance=.25cm}
            \diagarrow{from={1,3}, to={2,3}, text={i'}, x distance=.25cm}
            \diagarrow{from={1,1}, to={1,2}}
            \diagarrow{from={1,2}, to={1,3}}
            \diagarrow{from={2,1}, to={2,2}}
            \diagarrow{from={2,2}, to={2,3}}
        }}

        where the horizontal arrows compose to the identity.
        If $i\in M$ then $i'\in M$.

        \item $M$ is closed under countable many compositions and arbitrary direct sums.
        This means that if there is a countable chain
        $$ A_1 \xtos{i_1} A_2 \xtos{i_2} \cdots $$
        with all $i_j\in M$, then the canonical map $A_1\to\colim A_i$ is in $M$.

        Similarly if $i_j\colon A_j\to B_j$ for $j\in I$ are all $M$, then so is
        $$ \coprod_ji_j\colon\coprod_{j\in I}A_j \to \coprod_{j\in I}B_j . $$
    \eenum

\edefn

Note that the maps in (2), (3), (4) may not obviously be monic.
But these are all easily checked to be so.
We do this for (2), denote $f\colon A\to C$ so that $B\cup_AC$ is the disjoint union $B\amalg C$ where we identify
$f(a)$ and $i(a)$.
The map $i_*$ then maps $c\in C$ to its equivalence class $[c]$.
Notice that because $i$ is monic, we cannot identify $f(a),f(a')$ as that would mean $i(a)=i(a')$ and thus $a=a'$.
Thus no two distinct elements of $C$ are identified and thus $i_*$ is monic, as desired.

\bdefn

    A map $p\colon X\to Y$ is said to have the {\emphcolor right lifting property} ({\emphcolor RLP}) with respect to a class
    of morphisms $M$ if in any solid arrow diagram

    \centerline{\drawdiagram{
        $A$&$X$\cr
        $B$&$Y$
    }{
        \diagarrow{from={1,1}, to={1,2}}
        \diagarrow{from={1,1}, to={2,1}, text={i}, x distance=-.25cm}
        \diagarrow{from={2,1}, to={2,2}}
        \diagarrow{from={1,2}, to={2,2}, text={p}, x distance=.25cm}
        \diagarrow{from={2,1}, to={1,2}, dashed}
    }}

    where $i\in M$, a dashed arrow making the diagram commute exists.

    Dually we define the {\emphcolor left lifting property}.

\edefn

\blemm

    The class of monomorphisms $M_p$ having the left lifting property (LLP) with respect to a fixed map $p\colon X\to Y$
    is saturated.

\elemm

\bproof

    We prove condition (2).
    Suppose we have a commutative diagram

    \centerline{\drawdiagram{
        $A$&$C$&$X$\cr
        $B$&$B\cup_AC$&$Y$
    }{
        \diagarrow{from={1,1}, to={2,1}, text={i}, x distance=.25cm}
        \diagarrow{from={1,2}, to={2,2}, text={j}, x distance=.25cm}
        \diagarrow{from={1,3}, to={2,3}, text={p}, x distance=.25cm}
        \diagarrow{from={1,1}, to={1,2}}
        \diagarrow{from={1,2}, to={1,3}}
        \diagarrow{from={2,1}, to={2,2}}
        \diagarrow{from={2,1}, to={2,3}}
    }}

    where the left square is a pushout and $i\in M_p$; we wish to show $j\in M_p$.
    Because $i\in M_p$, there is a $\theta\colon B\to X$ so that

    \centerline{\drawdiagram{
        $A$&$X$\cr
        $B$&$Y$
    }{
        \diagarrow{from={1,1}, to={1,2}}
        \diagarrow{from={1,1}, to={2,1}, text={i}, x distance=-.25cm}
        \diagarrow{from={2,1}, to={2,2}}
        \diagarrow{from={1,2}, to={2,2}, text={p}, x distance=.25cm}
        \diagarrow{from={2,1}, to={1,2}, text={\theta}, x distance=-.25cm}
    }}

    commutes.
    But then this defines maps $B\to X,C\to X$ which commute and thus a map $\theta_*\colon B\cup_AC\to X$ which
    commutes with the necessary morphisms.
    Thus $j\in M_p$ as desired.
    \qed

\eproof

Clearly the intersection of saturated classes is saturated.
Thus for any class of monomorphisms $B$, we have a saturated class $M_B$ of all saturated classes containing $B$.
This is called the saturated class {\emphcolor generated} by $B$.
Alternatively $M_B$ is called the {\emphcolor saturation} of $B$.

Now consider the following classes of monomorphisms.
\bitems{1cm}
    \item{$\bB_1$} The set of all inclusions $\Lambda^n_k\subseteq\Delta^n$ for $0\leq k\leq n$, $n>0$.
    \item{$\bB_2$} The set of all inclusions
    $$ (\Delta^1\times\partial\Delta^n) \cup (\set e\times\Delta^n) \subseteq (\Delta^1\times\Delta^n),\qquad e=0,1 $$
    where $\set e$ corresponds to the simplicial set consisting of the vertex $e$ and all of its degeneracies.
    \item{$\bB_3$} The set of all inclusions
    $$ (\Delta^1\times Y) \cup (\set e\times X) \subseteq (\Delta^1\times X),\qquad e=0,1 $$
    where $Y\subseteq X$ is an inclusion.
\eitems

\bprop

    The classes $\bB_1$, $\bB_2$, and $\bB_3$ have the same saturation.

\eprop

\bproof

    Long and tedious.\CITE
    \qed

\eproof

\bdefn

    The saturation of $\bB_1$ (equivalently $\bB_2,\bB_3$) is called the class of {\emphcolor anodyne extensions}.

\edefn

\bcoro

    A fibration is a map $p\colon X\to Y$ which has the right lifting property with respect to all anodyne extensions.

\ecoro

\bproof

    Because the class of monics for which $p$ has the right lifting property is saturated, and contains
    $\Lambda^n_k\subseteq\Delta^n$ by definition, we obtain the desired result.
    \qed

\eproof

\bcoro[name={prodanodyne}]

    If $i\colon K\inj L$ is an anodyne extension, and $Y\inj X$ is an arbitrary inclusion, then
    $$ (K\times X)\cup(L\times Y) \to (L\times X) $$
    is an anodyne extension.

\ecoro

\bproof

    Consider the class of monics $K'\to L'$ such that
    $$ (K'\times X)\cup(L'\times Y) \to (L'\times X) $$
    is anodyne.
    This is a saturated set.

    Now consider the inclusion
    $$ (\Delta^1\times Y')\cup(\set e\times X')\subseteq(\Delta^1\times X'),\qquad Y'\subseteq X' $$
    We want to show that this is in the class, as it would imply all anodyne morphisms are.
    It is in the class iff the following inclusion is anodyne:
    $$ ((\Delta^1\times Y')\cup(\set e\times X'))\times X)\cup(\Delta^1\times X'\times Y) \to (\Delta^1\times X'\times X) $$
    which is just isomorphic to
    $$ (\Delta^1\times((Y'\times X)\cup(X'\times Y)))\cup(\set e\times X'\times X) \to \Delta^1\times X'\times X $$
    which is anodyne as desired.

    Thus all anodyne extensions are in the class, in particular $i$.
    \qed

\eproof

\bnote

    In the above corollary, the simplicial set $(K\times X)\cup(L\times Y)$ is the union when viewing both terms as subsets
    of $L\times X$.
    Categorically, this is the simplicial set $(K\times X)\cup_{K\times Y}(L\times Y)$.

\enote

